% Generated draft: requires mathematical and source review.
\documentclass[11pt]{article}
\usepackage[margin=25mm]{geometry}
\usepackage{fontspec}
\setmainfont{Latin Modern Roman}
\usepackage{amsmath,amssymb,mathtools,unicode-math}
\setmathfont{Latin Modern Math}
\usepackage{enumitem,needspace,xurl,hyperref}
\setlength{\parindent}{0pt}
\setlength{\parskip}{5pt}
\setlength{\emergencystretch}{3em}
\newcommand{\N}{\mathbb{N}}
\newcommand{\Z}{\mathbb{Z}}
\newcommand{\Q}{\mathbb{Q}}
\newcommand{\R}{\mathbb{R}}
\newcommand{\C}{\mathbb{C}}
\newcommand{\F}{\mathbb{F}}
\newcommand{\A}{\mathbb{A}}
\newcommand{\PP}{\mathbb{P}}
\newcommand{\E}{\mathbb{E}}
\newcommand{\eps}{\varepsilon}
\newcommand{\dd}{\,\mathrm{d}}
\begin{document}
\section*{UnsolvedMath ID: 30001257}
\textbf{Source problem code:} catalogue entry
\subsection{Definitions and mathematical statement}
The supplied record concerns blow-up in a sinh--Gordon problem. It does not specify the differential equation, its normalization, the underlying domain or manifold, boundary conditions, parameter regime, admissible solution families, or bounds imposed on those families. It also gives no definitions of a blow-up point or of its simple or multiple character. These missing specifications prevent a unique precise mathematical formulation.

Subject to recovering that setting and its terminology, let \(\mathcal B\) denote the blow-up points of an admissible solution family, and let \(\mathcal B_{\mathrm{s}}\) and \(\mathcal B_{\mathrm{m}}\) denote those called simple and multiple, respectively, under the intended definitions. These symbols are placeholders for definitions absent from the record; no exhaustiveness assertion about these classes is assumed.

The supplied classification question is: determine the possible simple or multiple character of the points in \(\mathcal B\), and characterize the circumstances in which
\[
\mathcal B_{\mathrm{m}}\ne\varnothing.
\]
In particular, the record does not determine whether multiplicity refers to several concentration profiles at one limiting point, clustering of concentration points, or several distinct blow-up locations. The displayed formulation therefore remains a conditional statement schema requiring recovery of the original hypotheses and definitions.
\subsection{Short English statement}
Classify the simple and multiple blow-up behavior in the intended sinh--Gordon problem, especially the circumstances permitting multiple blow-up. The supplied extract omits the equation, analytic hypotheses, and meaning of multiplicity, so a precise classification target remains unresolved.
\subsection{Sources}
\begin{itemize}
\item Supplied local UnsolvedMath registry record, permanent ID \(30001257\), source problem code \texttt{catalogue entry}, titled ``Multiplicity of Blow-Up Points in the Sinh--Gordon Equation.'' The record belongs to the Oberwolfach Reports Open Problems collection and gives the proposed year as 2009; no individual proposer is identified.
\item Source report: \emph{Topological and Variational Methods for Partial Differential Equations} (2009), Oberwolfach Reports, pp.\ 1427--1486, DOI \texttt{10.4171/owr/2009/26}. Supplied report link: \url{https://doi.org/10.4171/owr/2009/26}. No exact problem page or passage is supplied.
\item Supplied literature reference: Esposito--Wei, ``non-simple blow-up solutions,'' \url{https://doi.org/10.1007/s00526-008-0187-0}. This descriptive label is reproduced from the record and is not asserted to be the publication title.
\item Supplied literature reference: Jevnikar--Wei--Yang, ``classification of blow-up limits,'' \url{https://arxiv.org/abs/1602.02437}. This descriptive label is reproduced from the record and is not asserted to be the publication title.
\item Supplied literature assessment, checked on exactly \texttt{2026-08-21}: the record labels the problem partially solved, reports explicit non-simple blow-up solutions and later blow-up-value classifications, and states that a complete simple-versus-multiple classification in the extracted setting is not established. This is supplied dated curation derived from the OWR collection, not an independently verified literature assessment. The record expressly requests independent verification of time-sensitive status and cited claims.
\item The registry labels the statement \texttt{corrected\_verified}; its accompanying verification note says that the mathematical question was preserved while compound-name typography was normalized. This supplied label does not resolve the missing mathematical specifications.
\end{itemize}
\end{document}
